% =========================================================================
% preliminares/06_resumen.tex
% Resumen Ejecutivo y Palabras Clave
% =========================================================================
\clearpage
\thispagestyle{empty}

\begin{center}
  {\fontsize{14pt}{16pt}\selectfont\bfseries RESUMEN}
\end{center}
\vspace{0.5cm}

\fontsize{12pt}{14pt}\selectfont
El presente informe de prácticas preprofesionales describe el desarrollo y ejecución de actividades técnicas orientadas a resolver una problemática operativa en la institución receptora. En primer lugar, se describe la situación inicial y las limitaciones identificadas que motivaron la intervención. Posteriormente, se plantea el objetivo principal y los objetivos específicos orientados a diseñar, implementar y validar una solución robusta y estandarizada. Para alcanzar dichos propósitos, se aplicó una metodología de trabajo ágil estructurada en iteraciones periódicas, complementada con el uso de herramientas modernas de desarrollo de software, modelado de datos y aseguramiento de la calidad. Como resultado de las labores ejecutadas, se logró implementar con éxito los módulos funcionales requeridos, optimizando los flujos de trabajo internos y garantizando la estabilidad operativa del entorno. Finalmente, se concluye que la aplicación de principios de diseño y buenas prácticas de ingeniería permitió elevar significativamente la mantenibilidad, escalabilidad y eficiencia de los procesos en la organización.

\vspace{0.8cm}
\noindent\textbf{Palabras clave:} Prácticas preprofesionales, ingeniería de software, arquitectura de sistemas, optimización de procesos, desarrollo web.
\clearpage
